% ===== PRÉAMBULE CANONIQUE — à coller VERBATIM en tête de CHAQUE .tex =====
\documentclass[11pt,a4paper]{article}
\usepackage[utf8]{inputenc}\usepackage[T1]{fontenc}\usepackage{lmodern}
\usepackage[french]{babel}
\usepackage{amsmath,amssymb,mathtools}\usepackage{icomma}
\usepackage[version=4]{mhchem}          % \ce{...} : équations & formules
\usepackage{siunitx}\sisetup{locale=FR} % \qty{}{}, \si{}, \ang{}
\usepackage{chemfig}                    % \chemfig{...} : structures développées
\usepackage{xcolor}\usepackage{colortbl}
\usepackage{tikz}\usepackage{pgfplots}\pgfplotsset{compat=1.18}
\usepackage[most]{tcolorbox}\usepackage{enumitem}\usepackage{array,booktabs}
\usepackage[a4paper,margin=2.2cm]{geometry}
\usetikzlibrary{arrows.meta,calc,angles,quotes,positioning,patterns,backgrounds,shapes.geometric}
\definecolor{primary}{RGB}{222,49,99}\definecolor{primarydark}{RGB}{150,22,55}
\definecolor{defcol}{RGB}{0,114,178}\definecolor{methcol}{RGB}{52,92,148}
\definecolor{excol}{RGB}{0,135,90}\definecolor{attncol}{RGB}{200,40,55}
\tcbset{boxbase/.style={enhanced,breakable,boxrule=0.9pt,arc=2.5pt,
  left=3mm,right=3mm,top=2.3mm,bottom=2.3mm,fonttitle=\bfseries,coltitle=white}}
% --- boîtes TUTORIEL (noms EXACTS, ne pas renommer) ---
\newtcolorbox{cle}[1][]{boxbase,colback=defcol!6,colframe=defcol,
  title={\textsc{À retenir}\ifx\relax#1\relax\relax\else\textmd{ : \itshape #1}\fi}}
\newtcolorbox{methode}[1][]{boxbase,colback=methcol!7,colframe=methcol,sharp corners,
  title={\textsc{Méthode}\ifx\relax#1\relax\relax\else\textmd{ --- \itshape #1}\fi}}
\newtcolorbox{exemplebox}[1][]{boxbase,colback=excol!5,colframe=excol,
  title={\textsc{Exemple}\ifx\relax#1\relax\relax\else\textmd{ : \itshape #1}\fi}}
\newtcolorbox{piege}[1][]{boxbase,colback=attncol!6,colframe=attncol,
  title={\textsc{Piège}\ifx\relax#1\relax\relax\else\textmd{ : \itshape #1}\fi}}
\newtcolorbox{remarque}[1][]{boxbase,colback=black!4,colframe=black!45,
  title={\textsc{Remarque}\ifx\relax#1\relax\relax\else\textmd{ : \itshape #1}\fi}}
% --- boîtes EXERCICE / CORRIGÉ (noms EXACTS, auto-comptées) ---
\newtcolorbox[auto counter]{exo}[1][]{boxbase,colback=primary!4,colframe=primary,
  title={\textsc{Exercice}~\thetcbcounter\ifx\relax#1\relax\relax\else\textmd{ --- \itshape #1}\fi}}
\newtcolorbox[auto counter]{corrige}[1][]{boxbase,colback=excol!4,colframe=excol,
  title={\textsc{Corrigé}~\thetcbcounter\ifx\relax#1\relax\relax\else\textmd{ --- \itshape #1}\fi}}
\newcommand{\R}{\mathbb{R}}\newcommand{\N}{\mathbb{N}}\newcommand{\Z}{\mathbb{Z}}\newcommand{\ds}{\displaystyle}
% (un \usepackage{fancyhdr}+en-tête est possible ; il est ignoré côté web)
% =========================================================================
\begin{document}

\begin{corrige}[nitrate d'étain : lire la valence variable]
\begin{enumerate}[label=\alph*)]
  \item étain(IV) $=\ce{Sn^4+}$ (valence 4) + \ce{NO3-} (valence 1) $\Rightarrow$ quatre nitrates :
        \ce{Sn(NO3)4}. Vérif.\ : $+4$ et $4(-1)=-4$.
  \item étain(II) $=\ce{Sn^2+}$ (valence 2) + \ce{NO3-} $\Rightarrow$ \ce{Sn(NO3)2}.
  \item \ce{Sr} est le symbole du \emph{strontium}, pas de l'étain (\ce{Sn}). De plus \ce{Sr(NO3)2}
        correspondrait à un cation de valence 2, incompatible avec l'étain(IV).
\end{enumerate}
\end{corrige}

\begin{corrige}[une substance pure à partir de deux ions]
On croise \ce{Al^3+} (valence 3) avec l'anion de valence 2.
\begin{enumerate}[label=\alph*)]
  \item avec le sulfite \ce{SO3^2-} : \ce{Al2(SO3)3}.
  \item avec le sulfate \ce{SO4^2-} : \ce{Al2(SO4)3}.
  \item le sulf\emph{ite} \ce{SO3^2-} porte \emph{trois} oxygènes (soufre au degré $+\mathrm{IV}$), le
        sulf\emph{ate} \ce{SO4^2-} en porte \emph{quatre} (soufre au degré $+\mathrm{VI}$).
\end{enumerate}
\end{corrige}

\begin{corrige}[thiosulfate ou sulfite d'ammonium ?]
On croise \ce{NH4+} (valence 1) avec l'anion de valence 2 $\Rightarrow$ deux \ce{NH4+}, parenthésés.
\begin{enumerate}[label=\alph*)]
  \item thiosulfate d'ammonium : \ce{(NH4)2S2O3}.
  \item sulfite d'ammonium : \ce{(NH4)2SO3}.
  \item l'ion thiosulfate \ce{S2O3^2-} possède \emph{deux} atomes de soufre : c'est un sulfate dont un
        oxygène est remplacé par un soufre (préfixe \emph{thio-}). Le sulfite \ce{SO3^2-} n'a qu'un soufre.
\end{enumerate}
\end{corrige}

\begin{corrige}[identifier un métal alcalin par la masse molaire]
\begin{enumerate}[label=\alph*)]
  \item Un \emph{alcalin} est monovalent (\ce{M+}, valence 1) ; le phosphate \ce{PO4^3-} est de valence 3.
        Le croisement donne trois \ce{M+} pour un phosphate, d'où \ce{M3PO4}.
  \item $M_{\text{sel}}=\dfrac{m}{n}=\dfrac{42{,}4}{0{,}20}=\qty{212}{\gram\per\mol}$.
  \item $M(\ce{PO4})=31+4\times16=\qty{95}{\gram\per\mol}$, donc
        $3\,M(\ce{M})=212-95=117 \Rightarrow M(\ce{M})=\qty{39}{\gram\per\mol}$.
        C'est le \textbf{potassium} \ce{K} : le sel est \ce{K3PO4}.
\end{enumerate}
\end{corrige}

\begin{corrige}[retrouver un hydroxyde par sa composition massique]
\begin{enumerate}[label=\alph*)]
  \item Une masse molaire de \qty{23}{\gram\per\mol} correspond au \textbf{sodium} \ce{Na}.
  \item \ce{Na} est monovalent (\ce{Na+}) et \ce{OH-} est de valence 1 : $x=1$, $y=1$, soit \ce{NaOH}.
  \item $M(\ce{NaOH})=23+17=\qty{40}{\gram\per\mol}$, d'où
        $\dfrac{M(\ce{Na})}{M(\ce{NaOH})}=\dfrac{23}{40}=0{,}575=57{,}5\,\%$. Cohérent.
\end{enumerate}
\end{corrige}

\end{document}
